\chapter{Conjuntos}

\begin{objetivos}
  \item Introducir el lenguaje elemental de la teoría de conjuntos.
  \item Definir pertenencia, inclusión, unión, intersección, diferencia y complementario.
  \item Relacionar operaciones conjuntistas con conectores lógicos.
  \item Aprender a demostrar igualdades entre conjuntos.
\end{objetivos}

\section{Motivación}

La teoría de conjuntos proporciona uno de los lenguajes básicos de la matemática moderna. Muchas nociones que aparecerán a lo largo del curso, como aplicación, relación, estructura algebraica, número natural, número racional o espacio de soluciones, se describen de forma precisa mediante conjuntos y operaciones entre conjuntos.

En este capítulo no se pretende desarrollar una teoría axiomática de conjuntos. El objetivo es más elemental: introducir el vocabulario y la notación necesarios para trabajar con colecciones de objetos de forma rigurosa.

\begin{convencion}[Símbolos antes de su uso]
Todo símbolo conjuntista que aparezca en un enunciado debe haber sido definido previamente. En particular, antes de escribir expresiones como \(x\in A\), \(A\subseteq B\) o \(A\cup B\), se deberá haber indicado qué objetos representan \(x\), \(A\) y \(B\), y qué significado tienen los símbolos \(\in\), \(\subseteq\) y \(\cup\).
\end{convencion}

\section{Conjuntos y elementos}

De manera intuitiva, un \emph{conjunto} es una colección de objetos bien determinados. Los objetos que forman parte de un conjunto se llaman \emph{elementos} del conjunto.

\begin{definicion}[Conjunto y elemento]
Un \emph{conjunto} es una colección de objetos considerados como una totalidad. Si \(A\) es un conjunto, los objetos que pertenecen a \(A\) se llaman \emph{elementos} de \(A\).
\end{definicion}

La notación habitual consiste en denotar los conjuntos mediante letras mayúsculas, como \(A\), \(B\), \(C\), y sus elementos mediante letras minúsculas, como \(x\), \(y\), \(z\). Esta convención no es obligatoria, pero será utilizada de forma sistemática en estas notas.

\begin{ejemplo}
Sea \(A\) el conjunto formado por los números \(1\), \(2\) y \(3\). Escribimos
\[
  A=\{1,2,3\}.
\]
Los elementos de \(A\) son \(1\), \(2\) y \(3\).
\end{ejemplo}

Un conjunto puede describirse enumerando sus elementos, cuando esto es posible, o mediante una propiedad que caracteriza a sus elementos.

\begin{ejemplo}
El conjunto de los números naturales pares menores que \(10\) puede escribirse como
\[
  \{2,4,6,8\}.
\]
También puede describirse como
\[
  \{n\in\mathbb{N}: n \text{ es par y } n<10\}.
\]
En esta expresión, \(\mathbb{N}\) denota el conjunto de los números naturales y el símbolo ``:'' se lee como ``tal que''.
\end{ejemplo}

\begin{advertencia}
Para que una colección pueda considerarse un conjunto en el sentido elemental que utilizaremos en este curso, debe estar claro si un objeto dado pertenece o no pertenece a ella. Expresiones como ``el conjunto de los números grandes'' son ambiguas si no se precisa qué significa ``grande''.
\end{advertencia}

\section{Pertenencia e inclusión}

El primer símbolo fundamental de la teoría de conjuntos es el símbolo de pertenencia.

\begin{definicion}[Pertenencia]
Sean \(A\) un conjunto y \(x\) un objeto. Escribimos
\[
  x\in A
\]
para indicar que \(x\) es un elemento de \(A\). Escribimos
\[
  x\notin A
\]
para indicar que \(x\) no es un elemento de \(A\).
\end{definicion}

\begin{ejemplo}
Si \(A=\{1,2,3\}\), entonces
\[
  2\in A,
  \qquad
  5\notin A.
\]
\end{ejemplo}

El segundo símbolo fundamental es el símbolo de inclusión.

\begin{definicion}[Inclusión]
Sean \(A\) y \(B\) dos conjuntos. Decimos que \(A\) está \emph{incluido} en \(B\), o que \(A\) es un \emph{subconjunto} de \(B\), si todo elemento de \(A\) es también elemento de \(B\). En ese caso escribimos
\[
  A\subseteq B.
\]
En símbolos lógicos:
\[
  A\subseteq B
  \quad\Longleftrightarrow\quad
  \forall x\,(x\in A \Rightarrow x\in B).
\]
\end{definicion}

Aquí \(\forall\) significa ``para todo'', y \(\Rightarrow\) significa ``implica''. Es importante observar que estos símbolos se introdujeron en el capítulo anterior antes de ser utilizados aquí.

\begin{ejemplo}
Si \(A=\{1,2\}\) y \(B=\{1,2,3\}\), entonces \(A\subseteq B\), porque todo elemento de \(A\) es también elemento de \(B\).
\end{ejemplo}

\begin{advertencia}
No debe confundirse pertenencia con inclusión. La expresión \(x\in A\) afirma que \(x\) es un elemento del conjunto \(A\). La expresión \(A\subseteq B\) afirma que todo elemento de \(A\) es también elemento de \(B\). Por tanto, \(\in\) relaciona un objeto con un conjunto, mientras que \(\subseteq\) relaciona dos conjuntos.
\end{advertencia}

\section{Subconjuntos y conjunto vacío}

La relación de inclusión permite comparar conjuntos. Si \(A\subseteq B\), podemos decir que \(A\) es más pequeño o igual que \(B\) desde el punto de vista de sus elementos, aunque esta interpretación debe manejarse con cuidado en conjuntos infinitos.

\begin{definicion}[Igualdad de conjuntos]
Sean \(A\) y \(B\) dos conjuntos. Decimos que \(A\) y \(B\) son iguales, y escribimos \(A=B\), si tienen exactamente los mismos elementos. Equivalentemente,
\[
  A=B
  \quad\Longleftrightarrow\quad
  (A\subseteq B \text{ y } B\subseteq A).
\]
\end{definicion}

Esta caracterización es la base del método de \emph{doble inclusión}, que se estudiará al final del capítulo.

\begin{definicion}[Conjunto vacío]
El \emph{conjunto vacío} es el conjunto que no tiene ningún elemento. Se denota por
\[
  \varnothing.
\]
\end{definicion}

\begin{proposicion}
Para todo conjunto \(A\), se cumple
\[
  \varnothing\subseteq A.
\]
\end{proposicion}

\begin{proof}
Para probar que \(\varnothing\subseteq A\), debemos demostrar que todo elemento de \(\varnothing\) pertenece a \(A\). Pero \(\varnothing\) no tiene elementos. Por tanto, no existe ningún elemento de \(\varnothing\) que pueda violar la condición. En consecuencia, \(\varnothing\subseteq A\).
\end{proof}

\begin{advertencia}
El conjunto vacío \(\varnothing\) no debe confundirse con el conjunto \(\{\varnothing\}\). El primero no tiene elementos. El segundo tiene un elemento, que es precisamente \(\varnothing\).
\end{advertencia}

\section{Operaciones con conjuntos}

A partir de conjuntos dados se pueden construir nuevos conjuntos mediante operaciones. Las operaciones más importantes en este curso serán la unión, la intersección, la diferencia y el complementario.

En toda esta sección, sean \(A\) y \(B\) conjuntos contenidos en un conjunto de referencia \(U\), llamado \emph{universo} cuando sea necesario hablar de complementarios.

\subsection{Unión}

\begin{definicion}[Unión]
Sean \(A\) y \(B\) dos conjuntos. La \emph{unión} de \(A\) y \(B\) es el conjunto formado por los elementos que pertenecen a \(A\), a \(B\), o a ambos. Se denota por
\[
  A\cup B.
\]
Formalmente,
\[
  A\cup B = \{x : x\in A \text{ o } x\in B\}.
\]
\end{definicion}

La palabra ``o'' se interpreta aquí como una disyunción inclusiva: un elemento de \(A\cup B\) puede pertenecer a \(A\), a \(B\), o a ambos.

\begin{ejemplo}
Si \(A=\{1,2,3\}\) y \(B=\{3,4,5\}\), entonces
\[
  A\cup B=\{1,2,3,4,5\}.
\]
\end{ejemplo}

\subsection{Intersección}

\begin{definicion}[Intersección]
Sean \(A\) y \(B\) dos conjuntos. La \emph{intersección} de \(A\) y \(B\) es el conjunto formado por los elementos que pertenecen simultáneamente a \(A\) y a \(B\). Se denota por
\[
  A\cap B.
\]
Formalmente,
\[
  A\cap B = \{x : x\in A \text{ y } x\in B\}.
\]
\end{definicion}

\begin{ejemplo}
Si \(A=\{1,2,3\}\) y \(B=\{3,4,5\}\), entonces
\[
  A\cap B=\{3\}.
\]
\end{ejemplo}

\begin{definicion}[Conjuntos disjuntos]
Sean \(A\) y \(B\) dos conjuntos. Decimos que \(A\) y \(B\) son \emph{disjuntos} si no tienen elementos comunes, es decir, si
\[
  A\cap B=\varnothing.
\]
\end{definicion}

\subsection{Diferencia}

\begin{definicion}[Diferencia de conjuntos]
Sean \(A\) y \(B\) dos conjuntos. La \emph{diferencia} de \(A\) y \(B\) es el conjunto formado por los elementos que pertenecen a \(A\) y no pertenecen a \(B\). Se denota por
\[
  A\setminus B.
\]
Formalmente,
\[
  A\setminus B = \{x : x\in A \text{ y } x\notin B\}.
\]
\end{definicion}

\begin{ejemplo}
Si \(A=\{1,2,3\}\) y \(B=\{3,4,5\}\), entonces
\[
  A\setminus B=\{1,2\},
  \qquad
  B\setminus A=\{4,5\}.
\]
Por tanto, en general, \(A\setminus B\) no coincide con \(B\setminus A\).
\end{ejemplo}

\subsection{Complementario}

Para definir el complementario de un conjunto es necesario fijar previamente un conjunto de referencia.

\begin{definicion}[Complementario]
Sea \(U\) un conjunto de referencia y sea \(A\subseteq U\). El \emph{complementario} de \(A\) en \(U\) es el conjunto formado por los elementos de \(U\) que no pertenecen a \(A\). Se denota por
\[
  A^c
\]
cuando el conjunto de referencia \(U\) está claro. Formalmente,
\[
  A^c = \{x\in U : x\notin A\}.
\]
\end{definicion}

\begin{ejemplo}
Sea \(U=\{1,2,3,4,5\}\) y sea \(A=\{1,3,5\}\). Entonces
\[
  A^c=\{2,4\}.
\]
\end{ejemplo}

\begin{advertencia}
No tiene sentido hablar del complementario de \(A\) sin saber respecto de qué conjunto de referencia se toma el complementario. El conjunto \(A^c\) depende del universo \(U\).
\end{advertencia}

\section{Relación entre operaciones conjuntistas y conectores lógicos}

Las operaciones con conjuntos están estrechamente relacionadas con los conectores lógicos estudiados en el capítulo anterior.

Sean \(A\) y \(B\) subconjuntos de un conjunto de referencia \(U\), y sea \(x\in U\). Entonces:
\[
\begin{array}{rcl}
  x\in A\cup B &\Longleftrightarrow& x\in A \text{ o } x\in B,\\[2mm]
  x\in A\cap B &\Longleftrightarrow& x\in A \text{ y } x\in B,\\[2mm]
  x\in A\setminus B &\Longleftrightarrow& x\in A \text{ y } x\notin B,\\[2mm]
  x\in A^c &\Longleftrightarrow& x\notin A.
\end{array}
\]

Estas equivalencias permiten transformar problemas conjuntistas en problemas lógicos, y viceversa.

\begin{ejemplo}
La igualdad
\[
  (A\cap B)^c=A^c\cup B^c
\]
corresponde a la ley lógica
\[
  \neg(P\wedge Q) \Longleftrightarrow (\neg P)\vee(\neg Q),
\]
conocida como una de las leyes de De Morgan.
\end{ejemplo}

\section{Producto cartesiano}

Hasta ahora hemos trabajado con elementos individuales. En muchas situaciones matemáticas es necesario considerar pares ordenados.

\begin{definicion}[Par ordenado]
Un \emph{par ordenado} es una expresión de la forma \((a,b)\), donde \(a\) es la primera componente y \(b\) es la segunda componente. Dos pares ordenados \((a,b)\) y \((c,d)\) son iguales si y solo si
\[
  a=c \quad \text{y} \quad b=d.
\]
\end{definicion}

\begin{definicion}[Producto cartesiano]
Sean \(A\) y \(B\) dos conjuntos. El \emph{producto cartesiano} de \(A\) y \(B\) es el conjunto formado por todos los pares ordenados cuya primera componente pertenece a \(A\) y cuya segunda componente pertenece a \(B\). Se denota por
\[
  A\times B.
\]
Formalmente,
\[
  A\times B=\{(a,b): a\in A \text{ y } b\in B\}.
\]
\end{definicion}

\begin{ejemplo}
Si \(A=\{1,2\}\) y \(B=\{a,b\}\), entonces
\[
  A\times B=\{(1,a),(1,b),(2,a),(2,b)\}.
\]
\end{ejemplo}

\begin{advertencia}
En general, \(A\times B\) no coincide con \(B\times A\). Por ejemplo, si \(1\neq a\), entonces \((1,a)\) no es el mismo par ordenado que \((a,1)\).
\end{advertencia}

\section{Conjunto de las partes}

Además de considerar los elementos de un conjunto, también es importante considerar todos sus subconjuntos.

\begin{definicion}[Conjunto de las partes]
Sea \(A\) un conjunto. El \emph{conjunto de las partes} de \(A\), denotado por
\[
  \mathcal{P}(A),
\]
es el conjunto formado por todos los subconjuntos de \(A\). Formalmente,
\[
  \mathcal{P}(A)=\{B: B\subseteq A\}.
\]
\end{definicion}

\begin{ejemplo}
Si \(A=\{1,2\}\), entonces
\[
  \mathcal{P}(A)=\{\varnothing,\{1\},\{2\},\{1,2\}\}.
\]
\end{ejemplo}

\begin{proposicion}
Si \(A\) es un conjunto finito con \(n\) elementos, entonces \(\mathcal{P}(A)\) tiene \(2^n\) elementos.
\end{proposicion}

\begin{proof}
Cada subconjunto de \(A\) se obtiene decidiendo, para cada elemento de \(A\), si se incluye o no se incluye en el subconjunto. Para cada uno de los \(n\) elementos hay dos posibilidades. Por tanto, el número total de subconjuntos es
\[
  2\cdot 2\cdots 2 = 2^n.
\]
\end{proof}

\section{Familias de conjuntos}

En ocasiones se trabaja no solo con dos conjuntos, sino con muchos conjuntos a la vez. Para ello se usa la noción de familia de conjuntos.

\begin{definicion}[Familia de conjuntos]
Sea \(I\) un conjunto, llamado \emph{conjunto de índices}. Una \emph{familia de conjuntos indexada por} \(I\) es una colección de conjuntos
\[
  (A_i)_{i\in I},
\]
donde a cada índice \(i\in I\) le corresponde un conjunto \(A_i\).
\end{definicion}

\begin{ejemplo}
Si \(I=\{1,2,3\}\), una familia de conjuntos indexada por \(I\) puede escribirse como
\[
  (A_1,A_2,A_3).
\]
Por ejemplo,
\[
  A_1=\{1,2\},\qquad A_2=\{2,3\},\qquad A_3=\{3,4\}.
\]
\end{ejemplo}

\begin{definicion}[Unión e intersección de una familia]
Sea \((A_i)_{i\in I}\) una familia de conjuntos. La \emph{unión} de la familia se define por
\[
  \bigcup_{i\in I} A_i = \{x : \text{existe } i\in I \text{ tal que } x\in A_i\}.
\]
La \emph{intersección} de la familia se define por
\[
  \bigcap_{i\in I} A_i = \{x : \text{para todo } i\in I,\ x\in A_i\}.
\]
\end{definicion}

\begin{ejemplo}
Con los conjuntos
\[
  A_1=\{1,2\},\qquad A_2=\{2,3\},\qquad A_3=\{2,4\},
\]
se tiene
\[
  \bigcup_{i=1}^3 A_i = \{1,2,3,4\},
  \qquad
  \bigcap_{i=1}^3 A_i = \{2\}.
\]
\end{ejemplo}

\section{Demostración de identidades conjuntistas}

Una \emph{identidad conjuntista} es una igualdad entre conjuntos que se cumple bajo ciertas hipótesis. Por ejemplo,
\[
  A\cap(B\cup C)=(A\cap B)\cup(A\cap C)
\]
es una identidad conjuntista válida para cualesquiera conjuntos \(A\), \(B\) y \(C\).

El método más importante para demostrar igualdades entre conjuntos es el método de doble inclusión.

\begin{metodo}[Doble inclusión]
Para demostrar que dos conjuntos \(X\) e \(Y\) son iguales, basta demostrar las dos inclusiones:
\[
  X\subseteq Y
  \qquad\text{y}\qquad
  Y\subseteq X.
\]
Es decir, se prueba primero que todo elemento de \(X\) pertenece a \(Y\), y después que todo elemento de \(Y\) pertenece a \(X\).
\end{metodo}

\begin{proposicion}[Distributividad de la intersección respecto de la unión]
Sean \(A\), \(B\) y \(C\) conjuntos. Entonces
\[
  A\cap(B\cup C)=(A\cap B)\cup(A\cap C).
\]
\end{proposicion}

\begin{proof}
Demostraremos la igualdad mediante doble inclusión.

Primero probamos que
\[
  A\cap(B\cup C)\subseteq (A\cap B)\cup(A\cap C).
\]
Sea \(x\in A\cap(B\cup C)\). Entonces \(x\in A\) y \(x\in B\cup C\). Como \(x\in B\cup C\), se tiene \(x\in B\) o \(x\in C\).

Si \(x\in B\), entonces \(x\in A\cap B\), y por tanto
\[
  x\in (A\cap B)\cup(A\cap C).
\]
Si \(x\in C\), entonces \(x\in A\cap C\), y por tanto
\[
  x\in (A\cap B)\cup(A\cap C).
\]
En ambos casos, \(x\in (A\cap B)\cup(A\cap C)\). Por tanto,
\[
  A\cap(B\cup C)\subseteq (A\cap B)\cup(A\cap C).
\]

Ahora probamos la inclusión contraria:
\[
  (A\cap B)\cup(A\cap C)\subseteq A\cap(B\cup C).
\]
Sea \(x\in (A\cap B)\cup(A\cap C)\). Entonces \(x\in A\cap B\) o \(x\in A\cap C\).

Si \(x\in A\cap B\), entonces \(x\in A\) y \(x\in B\). En particular, \(x\in B\cup C\). Por tanto, \(x\in A\cap(B\cup C)\).

Si \(x\in A\cap C\), entonces \(x\in A\) y \(x\in C\). En particular, \(x\in B\cup C\). Por tanto, \(x\in A\cap(B\cup C)\).

En ambos casos, \(x\in A\cap(B\cup C)\). Por tanto,
\[
  (A\cap B)\cup(A\cap C)\subseteq A\cap(B\cup C).
\]

Como se cumplen ambas inclusiones, concluimos que
\[
  A\cap(B\cup C)=(A\cap B)\cup(A\cap C).
\]
\end{proof}

\begin{erroresfrecuentes}
  \item Confundir \(x\in A\) con \(x\subseteq A\).
  \item Confundir \(A\in B\) con \(A\subseteq B\).
  \item Escribir el complementario \(A^c\) sin haber fijado previamente el conjunto de referencia.
  \item Creer que \(A\setminus B = B\setminus A\) en general.
  \item Demostrar solo una inclusión y concluir indebidamente que dos conjuntos son iguales.
\end{erroresfrecuentes}

\begin{seminario}
  \item Demostración de igualdades conjuntistas mediante doble inclusión.
  \item Ejercicios de traducción entre lenguaje lógico y lenguaje conjuntista.
  \item Construcción de ejemplos finitos.
\end{seminario}

\begin{ejerciciospropuestos}
  \item Sean \(A=\{1,2,3\}\), \(B=\{2,3,4\}\) y \(C=\{3,4,5\}\). Calcula \(A\cup B\), \(A\cap B\), \(A\setminus B\), \(B\setminus A\), \(A\cap(B\cup C)\) y \((A\cap B)\cup(A\cap C)\).
  \item Demuestra, usando doble inclusión, la identidad
  \[
    A\cup(B\cap C)=(A\cup B)\cap(A\cup C).
  \]
  \item Demuestra las leyes de De Morgan:
  \[
    (A\cup B)^c=A^c\cap B^c,
    \qquad
    (A\cap B)^c=A^c\cup B^c.
  \]
  \item Calcula el conjunto de las partes de \(\{a,b\}\) y de \(\{a,b,c\}\).
  \item Sea \(A=\{1,2\}\) y \(B=\{x,y,z\}\). Escribe todos los elementos de \(A\times B\) y de \(B\times A\). ¿Son iguales estos dos productos cartesianos?
  \item Sea \((A_i)_{i\in\{1,2,3\}}\) la familia definida por
  \[
    A_1=\{1,2,3\},\qquad A_2=\{2,3,4\},\qquad A_3=\{3,4,5\}.
  \]
  Calcula \(\bigcup_{i=1}^3 A_i\) y \(\bigcap_{i=1}^3 A_i\).
\end{ejerciciospropuestos}

\resumen{El capítulo introduce el lenguaje elemental de la teoría de conjuntos. Las nociones de pertenencia, inclusión, operaciones conjuntistas, producto cartesiano, conjunto de las partes y familias de conjuntos proporcionan una base común para el resto del curso. La técnica central para demostrar igualdades entre conjuntos es la doble inclusión.}
